\documentclass[11pt,a4paper]{article}

\usepackage[utf8]{inputenc}
\usepackage[T1]{fontenc}
\usepackage{XCharter}
\usepackage[brazil]{babel}
\usepackage[margin=2cm]{geometry}
\usepackage{amsmath, amssymb}
\usepackage[xcharter,bigdelims,vvarbb]{newtxmath}
% newtxmath's operators font requests the fontspec family `XCharter(0)' in OT1;
% without these substitutions math digits and operator names silently fall back
% to Computer Modern (LaTeX Font Warning: Font shape `OT1/XCharter(0)/m/n' undefined).
\DeclareFontFamily{OT1}{XCharter(0)}{}
\DeclareFontShape{OT1}{XCharter(0)}{m}{n}{<->ssub * XCharter-TLF/m/n}{}
\DeclareFontShape{OT1}{XCharter(0)}{m}{it}{<->ssub * XCharter-TLF/m/it}{}
\DeclareFontShape{OT1}{XCharter(0)}{b}{n}{<->ssub * XCharter-TLF/b/n}{}
\DeclareFontShape{OT1}{XCharter(0)}{bx}{n}{<->ssub * XCharter-TLF/b/n}{}
\usepackage{enumerate}
\usepackage{array}
\usepackage{tikz}
\usetikzlibrary{positioning}

% ---- Painéis de função 1D usados na Questão 1 ---------------------------
% \fpA: eixo vertical em [-1, 1]; #1 coordenadas da poligonal, #2 rótulo.
% \fpB: eixo vertical em [-1,6; 1,6], usado nas alternativas; #1 coordenadas.
% Eixo horizontal sempre em [0, 2]; grade a cada 0,4 em x e a cada 0,2 em y.
\newcommand{\fpA}[2]{%
\begin{tikzpicture}[x=1.45cm, y=0.95cm, baseline=(current bounding box.center)]
  \foreach \gx in {0.4,0.8,1.2,1.6} \draw[gray!25] (\gx,-1) -- (\gx,1);
  \foreach \gy in {-0.8,-0.6,-0.4,-0.2,0.2,0.4,0.6,0.8} \draw[gray!20] (0,\gy) -- (2,\gy);
  \draw[gray!60, dashed] (0,0) -- (2,0);
  \draw[gray!70] (0,-1) rectangle (2,1);
  \draw[very thick, black!80] #1;
  \foreach \ty/\tl in {1/1,0.6/0{,}6,0.2/0{,}2,-0.2/-0{,}2,-0.6/-0{,}6,-1/-1}
    \node[font=\tiny, left, inner sep=1pt] at (0,\ty) {$\tl$};
  \foreach \tx/\tl in {0/0,0.4/0{,}4,0.8/0{,}8,1.2/1{,}2,1.6/1{,}6,2/2} \node[font=\tiny, below, inner sep=1.5pt] at (\tx,-1) {$\tl$};
  \node[font=\small, below, inner sep=1pt] at (1,-1.28) {#2};
\end{tikzpicture}}
\newcommand{\fpB}[1]{%
\begin{tikzpicture}[x=1.45cm, y=1.2cm, baseline=(current bounding box.center)]
  \foreach \gx in {0.4,0.8,1.2,1.6} \draw[gray!25] (\gx,-1.6) -- (\gx,1.6);
  \foreach \gy in {-1.4,-1.2,-1,-0.8,-0.6,-0.4,-0.2,0.2,0.4,0.6,0.8,1,1.2,1.4} \draw[gray!20] (0,\gy) -- (2,\gy);
  \draw[gray!60, dashed] (0,0) -- (2,0);
  \draw[gray!70] (0,-1.6) rectangle (2,1.6);
  \draw[very thick, black!80] #1;
  \foreach \ty/\tl in {1.6/1{,}6,1.2/1{,}2,0.8/0{,}8,0.4/0{,}4,0/0,-0.4/-0{,}4,-0.8/-0{,}8,-1.2/-1{,}2,-1.6/-1{,}6}
    \node[font=\tiny, left, inner sep=1pt] at (0,\ty) {$\tl$};
  \foreach \tx/\tl in {0/0,0.4/0{,}4,0.8/0{,}8,1.2/1{,}2,1.6/1{,}6,2/2} \node[font=\tiny, below, inner sep=1.5pt] at (\tx,-1.6) {$\tl$};
\end{tikzpicture}}

% ---- Histogramas da Questão 1 ------------------------------------------
% Pré-ativações de uma MLP ReLU simulada (n = 256, minilote de 1000 instâncias).
% Eixo horizontal comum em [-8, 8]; altura normalizada pelo pico de cada painel.
\newcommand{\histpanel}[1]{%
\begin{tikzpicture}[x=0.19cm, y=1.1cm, baseline=(current bounding box.center)]
  \foreach \gx in {-4,0,4} \draw[gray!25] (\gx,0) -- (\gx,1.05);
  \fill[black!22] plot[const plot] coordinates {#1} -- (8,0) -- (-8,0) -- cycle;
  \draw[black!75, thick] plot[const plot] coordinates {#1};
  \draw[gray!70] (-8,0) rectangle (8,1.05);
  \foreach \t in {-8,-4,0,4,8} \node[font=\tiny, below, inner sep=1.5pt] at (\t,0) {$\t$};
\end{tikzpicture}}

% ---- Mini-diagramas de rede (Questão 3) ---------------------------------
% Convenção: o viés fica dentro do círculo e os pesos ficam sobre as arestas.
\tikzset{
  mlpunit/.style={circle, draw, minimum size=8mm, inner sep=0pt, font=\small},
  mlpin/.style={font=\small},
  mlpw/.style={font=\scriptsize, inner sep=1.2pt}
}
% Região classificada como +1 (sombreada) no quadrado [0,4] x [0,4]; #1 é o polígono.
\newcommand{\regpanel}[1]{%
\begin{tikzpicture}[scale=0.64, baseline=(current bounding box.center)]
  \begin{scope}
    \clip (0,0) rectangle (4,4);
    \fill[black!18] #1;
  \end{scope}
  \foreach \g in {1,2,3} { \draw[gray!35] (\g,0) -- (\g,4); \draw[gray!35] (0,\g) -- (4,\g); }
  \draw[gray!70] (0,0) rectangle (4,4);
  \foreach \t in {0,1,2,3,4} { \node[font=\tiny, below, inner sep=1.5pt] at (\t,0) {$\t$}; \node[font=\tiny, left, inner sep=1.5pt] at (0,\t) {$\t$}; }
  \node[font=\scriptsize, below] at (2,-0.55) {$x_1$};
  \node[font=\scriptsize, left] at (-0.55,2) {$x_2$};
\end{tikzpicture}}

\pagestyle{empty}
\setlength{\parindent}{0pt}
\setlength{\parskip}{0.5em}

\renewcommand{\arraystretch}{1.4}

\begin{document}

%% =============================================================
%%  Header
%% =============================================================
{\Large\bfseries Aprendizado Profundo} \hfill \textbf{Prova, Unidade 1}\\
Prof. Lucas S. Kupssinskü \hfill Data: \rule{3cm}{0.4pt}\\[0.4em]
Nome: \rule{12cm}{0.4pt}

\rule{\linewidth}{0.8pt}

%% =============================================================
%%  Response grid
%% =============================================================
\textbf{Quadro de respostas (questões objetivas):}\\[0.2em]
\begin{center}
{\renewcommand{\arraystretch}{1.15}
\begin{tabular}{|c|c|c|c|c|c|c|c|c|}
\hline
\textbf{Questão} & \textbf{a} & \textbf{b} & \textbf{c} & \textbf{d} & \textbf{e} & \textbf{f} & \textbf{g} & \textbf{h} \\\hline
1 & & & & & & & & \\\hline
2 & & & & & & & & \\\hline
3 & & & & & & & & \\\hline
4 & & & & & & & & \\\hline
5 & & & & & & & & \\\hline
\end{tabular}
\hspace{1.2em}
\begin{tabular}{|c|c|c|c|c|c|c|c|c|}
\hline
\textbf{Questão} & \textbf{a} & \textbf{b} & \textbf{c} & \textbf{d} & \textbf{e} & \textbf{f} & \textbf{g} & \textbf{h} \\\hline
6 & & & & & & & & \\\hline
7 & & & & & & & & \\\hline
8 & & & & & & & & \\\hline
9 & & & & & & & & \\\hline
10 & & & & & & & & \\\hline
\end{tabular}}
\end{center}

\rule{\linewidth}{0.4pt}

%% =============================================================
%%  Objective questions  (10 x 1,0 = 10,0 pts)
%% =============================================================
\section*{Questões Objetivas (10{,}0 pts)}

\textit{Cada questão possui apenas uma alternativa correta. Marque a letra escolhida no quadro de respostas.}

\begin{enumerate}[1)]
\setlength{\itemsep}{1em}

%% ---------- Q1: OBJ-H inicialização, histogramas de pré-ativações em MLP ReLU ----------
\item (1{,}0) Um MLP tem entrada de dimensão $256$ e $4$ camadas ocultas com $n = 256$ unidades ReLU cada. As componentes da entrada são independentes, com distribuição $\mathcal{N}(0, 1)$, os vieses são inicializados em zero e os pesos com $\theta^{(l)}_{ij} \sim \mathcal{N}(0, \sigma_\theta^2)$. A rede foi inicializada com quatro valores diferentes de $\sigma_\theta^2$ e, para cada um, mediu-se o histograma das pré-ativações $\mathbf{Z}^{(1)}, \ldots, \mathbf{Z}^{(4)}$ em um minilote de $1000$ instâncias. Cada linha (P, Q, R e S) corresponde a uma inicialização. O eixo horizontal é o mesmo em todos os painéis, e a altura de cada painel está normalizada pelo seu próprio pico.

\begin{center}
\setlength{\tabcolsep}{3pt}
\begin{tabular}{c@{\hspace{6pt}}cccc}
 & $\mathbf{Z}^{(1)}$ & $\mathbf{Z}^{(2)}$ & $\mathbf{Z}^{(3)}$ & $\mathbf{Z}^{(4)}$ \\
\textbf{P} & \histpanel{(-8.0,0.0000) (-7.8,0.0000) (-7.6,0.0000) (-7.4,0.0000) (-7.2,0.0000) (-7.0,0.0000) (-6.8,0.0000) (-6.6,0.0000) (-6.4,0.0000) (-6.2,0.0000) (-6.0,0.0000) (-5.8,0.0000) (-5.6,0.0000) (-5.4,0.0000) (-5.2,0.0000) (-5.0,0.0000) (-4.8,0.0000) (-4.6,0.0000) (-4.4,0.0002) (-4.2,0.0003) (-4.0,0.0004) (-3.8,0.0010) (-3.6,0.0025) (-3.4,0.0049) (-3.2,0.0080) (-3.0,0.0154) (-2.8,0.0286) (-2.6,0.0483) (-2.4,0.0750) (-2.2,0.1118) (-2.0,0.1662) (-1.8,0.2334) (-1.6,0.3302) (-1.4,0.4315) (-1.2,0.5589) (-1.0,0.6683) (-0.8,0.7911) (-0.6,0.8984) (-0.4,0.9539) (-0.2,1.0000) (0.0,0.9889) (0.2,0.9646) (0.4,0.8935) (0.6,0.7910) (0.8,0.6685) (1.0,0.5409) (1.2,0.4260) (1.4,0.3307) (1.6,0.2295) (1.8,0.1693) (2.0,0.1123) (2.2,0.0742) (2.4,0.0458) (2.6,0.0265) (2.8,0.0168) (3.0,0.0083) (3.2,0.0053) (3.4,0.0027) (3.6,0.0014) (3.8,0.0008) (4.0,0.0002) (4.2,0.0000) (4.4,0.0001) (4.6,0.0000) (4.8,0.0000) (5.0,0.0000) (5.2,0.0000) (5.4,0.0000) (5.6,0.0000) (5.8,0.0000) (6.0,0.0000) (6.2,0.0000) (6.4,0.0000) (6.6,0.0000) (6.8,0.0000) (7.0,0.0000) (7.2,0.0000) (7.4,0.0000) (7.6,0.0000) (7.8,0.0000) (8.0,0)} & \histpanel{(-8.0,0.0000) (-7.8,0.0000) (-7.6,0.0000) (-7.4,0.0000) (-7.2,0.0000) (-7.0,0.0000) (-6.8,0.0000) (-6.6,0.0000) (-6.4,0.0000) (-6.2,0.0000) (-6.0,0.0000) (-5.8,0.0000) (-5.6,0.0000) (-5.4,0.0000) (-5.2,0.0000) (-5.0,0.0000) (-4.8,0.0000) (-4.6,0.0000) (-4.4,0.0000) (-4.2,0.0000) (-4.0,0.0000) (-3.8,0.0000) (-3.6,0.0000) (-3.4,0.0001) (-3.2,0.0001) (-3.0,0.0004) (-2.8,0.0011) (-2.6,0.0029) (-2.4,0.0066) (-2.2,0.0141) (-2.0,0.0313) (-1.8,0.0622) (-1.6,0.1152) (-1.4,0.1978) (-1.2,0.3102) (-1.0,0.4550) (-0.8,0.6305) (-0.6,0.8029) (-0.4,0.9377) (-0.2,1.0000) (0.0,0.9893) (0.2,0.8887) (0.4,0.7502) (0.6,0.5812) (0.8,0.4017) (1.0,0.2626) (1.2,0.1576) (1.4,0.0863) (1.6,0.0452) (1.8,0.0214) (2.0,0.0095) (2.2,0.0040) (2.4,0.0019) (2.6,0.0005) (2.8,0.0004) (3.0,0.0002) (3.2,0.0000) (3.4,0.0000) (3.6,0.0000) (3.8,0.0000) (4.0,0.0000) (4.2,0.0000) (4.4,0.0000) (4.6,0.0000) (4.8,0.0000) (5.0,0.0000) (5.2,0.0000) (5.4,0.0000) (5.6,0.0000) (5.8,0.0000) (6.0,0.0000) (6.2,0.0000) (6.4,0.0000) (6.6,0.0000) (6.8,0.0000) (7.0,0.0000) (7.2,0.0000) (7.4,0.0000) (7.6,0.0000) (7.8,0.0000) (8.0,0)} & \histpanel{(-8.0,0.0000) (-7.8,0.0000) (-7.6,0.0000) (-7.4,0.0000) (-7.2,0.0000) (-7.0,0.0000) (-6.8,0.0000) (-6.6,0.0000) (-6.4,0.0000) (-6.2,0.0000) (-6.0,0.0000) (-5.8,0.0000) (-5.6,0.0000) (-5.4,0.0000) (-5.2,0.0000) (-5.0,0.0000) (-4.8,0.0000) (-4.6,0.0000) (-4.4,0.0000) (-4.2,0.0000) (-4.0,0.0000) (-3.8,0.0000) (-3.6,0.0000) (-3.4,0.0000) (-3.2,0.0000) (-3.0,0.0000) (-2.8,0.0000) (-2.6,0.0000) (-2.4,0.0000) (-2.2,0.0002) (-2.0,0.0006) (-1.8,0.0023) (-1.6,0.0072) (-1.4,0.0202) (-1.2,0.0549) (-1.0,0.1355) (-0.8,0.2918) (-0.6,0.5443) (-0.4,0.8085) (-0.2,1.0000) (0.0,0.9973) (0.2,0.8191) (0.4,0.5652) (0.6,0.3119) (0.8,0.1454) (1.0,0.0596) (1.2,0.0214) (1.4,0.0075) (1.6,0.0018) (1.8,0.0004) (2.0,0.0001) (2.2,0.0000) (2.4,0.0000) (2.6,0.0000) (2.8,0.0000) (3.0,0.0000) (3.2,0.0000) (3.4,0.0000) (3.6,0.0000) (3.8,0.0000) (4.0,0.0000) (4.2,0.0000) (4.4,0.0000) (4.6,0.0000) (4.8,0.0000) (5.0,0.0000) (5.2,0.0000) (5.4,0.0000) (5.6,0.0000) (5.8,0.0000) (6.0,0.0000) (6.2,0.0000) (6.4,0.0000) (6.6,0.0000) (6.8,0.0000) (7.0,0.0000) (7.2,0.0000) (7.4,0.0000) (7.6,0.0000) (7.8,0.0000) (8.0,0)} & \histpanel{(-8.0,0.0000) (-7.8,0.0000) (-7.6,0.0000) (-7.4,0.0000) (-7.2,0.0000) (-7.0,0.0000) (-6.8,0.0000) (-6.6,0.0000) (-6.4,0.0000) (-6.2,0.0000) (-6.0,0.0000) (-5.8,0.0000) (-5.6,0.0000) (-5.4,0.0000) (-5.2,0.0000) (-5.0,0.0000) (-4.8,0.0000) (-4.6,0.0000) (-4.4,0.0000) (-4.2,0.0000) (-4.0,0.0000) (-3.8,0.0000) (-3.6,0.0000) (-3.4,0.0000) (-3.2,0.0000) (-3.0,0.0000) (-2.8,0.0000) (-2.6,0.0000) (-2.4,0.0000) (-2.2,0.0000) (-2.0,0.0000) (-1.8,0.0000) (-1.6,0.0001) (-1.4,0.0013) (-1.2,0.0085) (-1.0,0.0403) (-0.8,0.1341) (-0.6,0.3412) (-0.4,0.6837) (-0.2,1.0000) (0.0,0.9993) (0.2,0.6824) (0.4,0.3201) (0.6,0.1117) (0.8,0.0283) (1.0,0.0054) (1.2,0.0006) (1.4,0.0001) (1.6,0.0000) (1.8,0.0000) (2.0,0.0000) (2.2,0.0000) (2.4,0.0000) (2.6,0.0000) (2.8,0.0000) (3.0,0.0000) (3.2,0.0000) (3.4,0.0000) (3.6,0.0000) (3.8,0.0000) (4.0,0.0000) (4.2,0.0000) (4.4,0.0000) (4.6,0.0000) (4.8,0.0000) (5.0,0.0000) (5.2,0.0000) (5.4,0.0000) (5.6,0.0000) (5.8,0.0000) (6.0,0.0000) (6.2,0.0000) (6.4,0.0000) (6.6,0.0000) (6.8,0.0000) (7.0,0.0000) (7.2,0.0000) (7.4,0.0000) (7.6,0.0000) (7.8,0.0000) (8.0,0)} \\[0.4em]
\textbf{Q} & \histpanel{(-8.0,0.0004) (-7.8,0.0008) (-7.6,0.0003) (-7.4,0.0010) (-7.2,0.0021) (-7.0,0.0027) (-6.8,0.0041) (-6.6,0.0040) (-6.4,0.0073) (-6.2,0.0103) (-6.0,0.0125) (-5.8,0.0162) (-5.6,0.0228) (-5.4,0.0282) (-5.2,0.0388) (-5.0,0.0479) (-4.8,0.0597) (-4.6,0.0763) (-4.4,0.0961) (-4.2,0.1188) (-4.0,0.1496) (-3.8,0.1776) (-3.6,0.2132) (-3.4,0.2452) (-3.2,0.2967) (-3.0,0.3533) (-2.8,0.4026) (-2.6,0.4559) (-2.4,0.5169) (-2.2,0.5604) (-2.0,0.6292) (-1.8,0.6916) (-1.6,0.7478) (-1.4,0.8241) (-1.2,0.8594) (-1.0,0.9059) (-0.8,0.9593) (-0.6,0.9930) (-0.4,0.9950) (-0.2,0.9888) (0.0,1.0000) (0.2,0.9924) (0.4,0.9851) (0.6,0.9430) (0.8,0.9168) (1.0,0.8490) (1.2,0.8132) (1.4,0.7601) (1.6,0.6936) (1.8,0.6366) (2.0,0.5755) (2.2,0.5081) (2.4,0.4543) (2.6,0.3941) (2.8,0.3440) (3.0,0.2995) (3.2,0.2577) (3.4,0.2085) (3.6,0.1794) (3.8,0.1485) (4.0,0.1214) (4.2,0.1030) (4.4,0.0707) (4.6,0.0635) (4.8,0.0471) (5.0,0.0420) (5.2,0.0282) (5.4,0.0250) (5.6,0.0179) (5.8,0.0131) (6.0,0.0106) (6.2,0.0073) (6.4,0.0039) (6.6,0.0031) (6.8,0.0027) (7.0,0.0018) (7.2,0.0016) (7.4,0.0007) (7.6,0.0010) (7.8,0.0001) (8.0,0)} & \histpanel{(-8.0,0.0207) (-7.8,0.0277) (-7.6,0.0260) (-7.4,0.0380) (-7.2,0.0399) (-7.0,0.0467) (-6.8,0.0588) (-6.6,0.0644) (-6.4,0.0817) (-6.2,0.0909) (-6.0,0.0999) (-5.8,0.1206) (-5.6,0.1377) (-5.4,0.1550) (-5.2,0.1718) (-5.0,0.1984) (-4.8,0.2264) (-4.6,0.2550) (-4.4,0.2801) (-4.2,0.3219) (-4.0,0.3681) (-3.8,0.3852) (-3.6,0.4293) (-3.4,0.4774) (-3.2,0.5005) (-3.0,0.5459) (-2.8,0.5835) (-2.6,0.6350) (-2.4,0.6702) (-2.2,0.7243) (-2.0,0.7582) (-1.8,0.8047) (-1.6,0.8490) (-1.4,0.8652) (-1.2,0.8985) (-1.0,0.9313) (-0.8,0.9394) (-0.6,0.9567) (-0.4,0.9827) (-0.2,0.9849) (0.0,1.0000) (0.2,0.9705) (0.4,0.9684) (0.6,0.9433) (0.8,0.9262) (1.0,0.9102) (1.2,0.8848) (1.4,0.8491) (1.6,0.8167) (1.8,0.7732) (2.0,0.7167) (2.2,0.7045) (2.4,0.6554) (2.6,0.6096) (2.8,0.5750) (3.0,0.5166) (3.2,0.4954) (3.4,0.4493) (3.6,0.4027) (3.8,0.3684) (4.0,0.3333) (4.2,0.3040) (4.4,0.2593) (4.6,0.2399) (4.8,0.2049) (5.0,0.1756) (5.2,0.1607) (5.4,0.1454) (5.6,0.1121) (5.8,0.1101) (6.0,0.0876) (6.2,0.0730) (6.4,0.0675) (6.6,0.0523) (6.8,0.0467) (7.0,0.0393) (7.2,0.0316) (7.4,0.0309) (7.6,0.0216) (7.8,0.0190) (8.0,0)} & \histpanel{(-8.0,0.1451) (-7.8,0.1544) (-7.6,0.1620) (-7.4,0.1762) (-7.2,0.2013) (-7.0,0.2163) (-6.8,0.2403) (-6.6,0.2580) (-6.4,0.2833) (-6.2,0.2918) (-6.0,0.3051) (-5.8,0.3530) (-5.6,0.3724) (-5.4,0.4030) (-5.2,0.4234) (-5.0,0.4608) (-4.8,0.4964) (-4.6,0.5131) (-4.4,0.5342) (-4.2,0.5549) (-4.0,0.6025) (-3.8,0.6337) (-3.6,0.6601) (-3.4,0.6894) (-3.2,0.7196) (-3.0,0.7433) (-2.8,0.7587) (-2.6,0.8266) (-2.4,0.8209) (-2.2,0.8515) (-2.0,0.8624) (-1.8,0.9040) (-1.6,0.8985) (-1.4,0.9348) (-1.2,0.9570) (-1.0,0.9618) (-0.8,0.9800) (-0.6,0.9641) (-0.4,1.0000) (-0.2,0.9842) (0.0,0.9740) (0.2,0.9643) (0.4,0.9532) (0.6,0.9810) (0.8,0.9601) (1.0,0.9382) (1.2,0.9363) (1.4,0.9038) (1.6,0.8899) (1.8,0.8568) (2.0,0.8260) (2.2,0.8101) (2.4,0.8011) (2.6,0.7656) (2.8,0.7477) (3.0,0.6966) (3.2,0.6865) (3.4,0.6667) (3.6,0.6468) (3.8,0.5924) (4.0,0.5709) (4.2,0.5422) (4.4,0.5171) (4.6,0.4776) (4.8,0.4479) (5.0,0.4175) (5.2,0.3726) (5.4,0.3719) (5.6,0.3348) (5.8,0.3072) (6.0,0.2913) (6.2,0.2766) (6.4,0.2439) (6.6,0.2340) (6.8,0.2169) (7.0,0.2063) (7.2,0.1768) (7.4,0.1675) (7.6,0.1544) (7.8,0.1308) (8.0,0)} & \histpanel{(-8.0,0.2992) (-7.8,0.3500) (-7.6,0.3401) (-7.4,0.3920) (-7.2,0.3887) (-7.0,0.4218) (-6.8,0.4205) (-6.6,0.4345) (-6.4,0.4744) (-6.2,0.5075) (-6.0,0.5088) (-5.8,0.5182) (-5.6,0.5630) (-5.4,0.5825) (-5.2,0.6026) (-5.0,0.6258) (-4.8,0.6372) (-4.6,0.6660) (-4.4,0.7192) (-4.2,0.7024) (-4.0,0.7278) (-3.8,0.7235) (-3.6,0.7777) (-3.4,0.7937) (-3.2,0.7856) (-3.0,0.8067) (-2.8,0.8166) (-2.6,0.8731) (-2.4,0.8649) (-2.2,0.9125) (-2.0,0.9361) (-1.8,0.9191) (-1.6,0.9570) (-1.4,0.9423) (-1.2,0.9316) (-1.0,0.9415) (-0.8,0.9679) (-0.6,0.9819) (-0.4,0.9791) (-0.2,1.0000) (0.0,0.9763) (0.2,0.9845) (0.4,0.9921) (0.6,0.9789) (0.8,0.9524) (1.0,0.9506) (1.2,0.9359) (1.4,0.9428) (1.6,0.9303) (1.8,0.9387) (2.0,0.8949) (2.2,0.8692) (2.4,0.8603) (2.6,0.8408) (2.8,0.8293) (3.0,0.8123) (3.2,0.7833) (3.4,0.7738) (3.6,0.7693) (3.8,0.7548) (4.0,0.7108) (4.2,0.6909) (4.4,0.6467) (4.6,0.6383) (4.8,0.6324) (5.0,0.5714) (5.2,0.5843) (5.4,0.5495) (5.6,0.5429) (5.8,0.5162) (6.0,0.4765) (6.2,0.4864) (6.4,0.4582) (6.6,0.4172) (6.8,0.3984) (7.0,0.3773) (7.2,0.3625) (7.4,0.3661) (7.6,0.3284) (7.8,0.2989) (8.0,0)} \\[0.4em]
\textbf{R} & \histpanel{(-8.0,0.0000) (-7.8,0.0000) (-7.6,0.0000) (-7.4,0.0000) (-7.2,0.0000) (-7.0,0.0000) (-6.8,0.0000) (-6.6,0.0000) (-6.4,0.0001) (-6.2,0.0000) (-6.0,0.0002) (-5.8,0.0001) (-5.6,0.0010) (-5.4,0.0010) (-5.2,0.0017) (-5.0,0.0025) (-4.8,0.0043) (-4.6,0.0060) (-4.4,0.0089) (-4.2,0.0151) (-4.0,0.0237) (-3.8,0.0312) (-3.6,0.0445) (-3.4,0.0629) (-3.2,0.0892) (-3.0,0.1196) (-2.8,0.1591) (-2.6,0.2022) (-2.4,0.2625) (-2.2,0.3245) (-2.0,0.3977) (-1.8,0.4713) (-1.6,0.5594) (-1.4,0.6373) (-1.2,0.7382) (-1.0,0.8258) (-0.8,0.8843) (-0.6,0.9382) (-0.4,0.9670) (-0.2,1.0000) (0.0,0.9905) (0.2,0.9637) (0.4,0.9250) (0.6,0.8881) (0.8,0.8026) (1.0,0.7248) (1.2,0.6562) (1.4,0.5501) (1.6,0.4810) (1.8,0.4096) (2.0,0.3252) (2.2,0.2673) (2.4,0.2070) (2.6,0.1620) (2.8,0.1211) (3.0,0.0890) (3.2,0.0620) (3.4,0.0454) (3.6,0.0343) (3.8,0.0230) (4.0,0.0140) (4.2,0.0108) (4.4,0.0053) (4.6,0.0049) (4.8,0.0017) (5.0,0.0020) (5.2,0.0011) (5.4,0.0010) (5.6,0.0003) (5.8,0.0001) (6.0,0.0001) (6.2,0.0001) (6.4,0.0000) (6.6,0.0000) (6.8,0.0000) (7.0,0.0000) (7.2,0.0000) (7.4,0.0000) (7.6,0.0000) (7.8,0.0000) (8.0,0)} & \histpanel{(-8.0,0.0000) (-7.8,0.0000) (-7.6,0.0000) (-7.4,0.0000) (-7.2,0.0000) (-7.0,0.0001) (-6.8,0.0000) (-6.6,0.0000) (-6.4,0.0003) (-6.2,0.0002) (-6.0,0.0004) (-5.8,0.0004) (-5.6,0.0010) (-5.4,0.0013) (-5.2,0.0024) (-5.0,0.0047) (-4.8,0.0075) (-4.6,0.0102) (-4.4,0.0135) (-4.2,0.0228) (-4.0,0.0314) (-3.8,0.0413) (-3.6,0.0600) (-3.4,0.0777) (-3.2,0.1059) (-3.0,0.1442) (-2.8,0.1880) (-2.6,0.2438) (-2.4,0.2939) (-2.2,0.3605) (-2.0,0.4287) (-1.8,0.5107) (-1.6,0.5769) (-1.4,0.6748) (-1.2,0.7465) (-1.0,0.8360) (-0.8,0.8968) (-0.6,0.9282) (-0.4,0.9790) (-0.2,1.0000) (0.0,0.9960) (0.2,0.9712) (0.4,0.9303) (0.6,0.8618) (0.8,0.8161) (1.0,0.7234) (1.2,0.6382) (1.4,0.5546) (1.6,0.4737) (1.8,0.4029) (2.0,0.3260) (2.2,0.2728) (2.4,0.2051) (2.6,0.1600) (2.8,0.1209) (3.0,0.0894) (3.2,0.0679) (3.4,0.0475) (3.6,0.0336) (3.8,0.0219) (4.0,0.0172) (4.2,0.0118) (4.4,0.0086) (4.6,0.0052) (4.8,0.0033) (5.0,0.0027) (5.2,0.0012) (5.4,0.0011) (5.6,0.0006) (5.8,0.0005) (6.0,0.0001) (6.2,0.0001) (6.4,0.0000) (6.6,0.0000) (6.8,0.0001) (7.0,0.0001) (7.2,0.0000) (7.4,0.0000) (7.6,0.0000) (7.8,0.0000) (8.0,0)} & \histpanel{(-8.0,0.0000) (-7.8,0.0000) (-7.6,0.0000) (-7.4,0.0001) (-7.2,0.0000) (-7.0,0.0000) (-6.8,0.0001) (-6.6,0.0000) (-6.4,0.0003) (-6.2,0.0002) (-6.0,0.0008) (-5.8,0.0010) (-5.6,0.0013) (-5.4,0.0030) (-5.2,0.0037) (-5.0,0.0050) (-4.8,0.0074) (-4.6,0.0110) (-4.4,0.0164) (-4.2,0.0202) (-4.0,0.0320) (-3.8,0.0410) (-3.6,0.0574) (-3.4,0.0756) (-3.2,0.1030) (-3.0,0.1339) (-2.8,0.1800) (-2.6,0.2261) (-2.4,0.2870) (-2.2,0.3525) (-2.0,0.4241) (-1.8,0.4993) (-1.6,0.5979) (-1.4,0.6756) (-1.2,0.7714) (-1.0,0.8382) (-0.8,0.9103) (-0.6,0.9438) (-0.4,0.9863) (-0.2,1.0000) (0.0,0.9916) (0.2,0.9533) (0.4,0.9271) (0.6,0.8647) (0.8,0.8161) (1.0,0.7314) (1.2,0.6583) (1.4,0.5683) (1.6,0.4868) (1.8,0.4253) (2.0,0.3398) (2.2,0.2852) (2.4,0.2178) (2.6,0.1741) (2.8,0.1366) (3.0,0.1041) (3.2,0.0740) (3.4,0.0531) (3.6,0.0397) (3.8,0.0270) (4.0,0.0176) (4.2,0.0118) (4.4,0.0082) (4.6,0.0052) (4.8,0.0030) (5.0,0.0019) (5.2,0.0011) (5.4,0.0007) (5.6,0.0004) (5.8,0.0001) (6.0,0.0001) (6.2,0.0001) (6.4,0.0000) (6.6,0.0000) (6.8,0.0000) (7.0,0.0000) (7.2,0.0000) (7.4,0.0000) (7.6,0.0000) (7.8,0.0000) (8.0,0)} & \histpanel{(-8.0,0.0000) (-7.8,0.0000) (-7.6,0.0000) (-7.4,0.0001) (-7.2,0.0001) (-7.0,0.0000) (-6.8,0.0001) (-6.6,0.0001) (-6.4,0.0003) (-6.2,0.0003) (-6.0,0.0011) (-5.8,0.0012) (-5.6,0.0010) (-5.4,0.0019) (-5.2,0.0030) (-5.0,0.0042) (-4.8,0.0062) (-4.6,0.0101) (-4.4,0.0136) (-4.2,0.0218) (-4.0,0.0266) (-3.8,0.0379) (-3.6,0.0523) (-3.4,0.0698) (-3.2,0.0930) (-3.0,0.1230) (-2.8,0.1558) (-2.6,0.1977) (-2.4,0.2637) (-2.2,0.3338) (-2.0,0.4050) (-1.8,0.4861) (-1.6,0.5879) (-1.4,0.6653) (-1.2,0.7569) (-1.0,0.8493) (-0.8,0.9285) (-0.6,0.9612) (-0.4,0.9932) (-0.2,1.0000) (0.0,0.9963) (0.2,0.9660) (0.4,0.9234) (0.6,0.8554) (0.8,0.7935) (1.0,0.7182) (1.2,0.6268) (1.4,0.5377) (1.6,0.4629) (1.8,0.3834) (2.0,0.3139) (2.2,0.2556) (2.4,0.1982) (2.6,0.1563) (2.8,0.1208) (3.0,0.0938) (3.2,0.0684) (3.4,0.0498) (3.6,0.0375) (3.8,0.0263) (4.0,0.0182) (4.2,0.0126) (4.4,0.0095) (4.6,0.0066) (4.8,0.0045) (5.0,0.0024) (5.2,0.0011) (5.4,0.0014) (5.6,0.0004) (5.8,0.0005) (6.0,0.0001) (6.2,0.0001) (6.4,0.0001) (6.6,0.0001) (6.8,0.0000) (7.0,0.0000) (7.2,0.0000) (7.4,0.0000) (7.6,0.0000) (7.8,0.0000) (8.0,0)} \\[0.4em]
\textbf{S} & \histpanel{(-8.0,0.0000) (-7.8,0.0000) (-7.6,0.0000) (-7.4,0.0000) (-7.2,0.0000) (-7.0,0.0000) (-6.8,0.0000) (-6.6,0.0000) (-6.4,0.0000) (-6.2,0.0000) (-6.0,0.0000) (-5.8,0.0000) (-5.6,0.0000) (-5.4,0.0000) (-5.2,0.0000) (-5.0,0.0000) (-4.8,0.0000) (-4.6,0.0000) (-4.4,0.0000) (-4.2,0.0000) (-4.0,0.0000) (-3.8,0.0000) (-3.6,0.0000) (-3.4,0.0000) (-3.2,0.0000) (-3.0,0.0000) (-2.8,0.0000) (-2.6,0.0000) (-2.4,0.0000) (-2.2,0.0001) (-2.0,0.0009) (-1.8,0.0036) (-1.6,0.0119) (-1.4,0.0353) (-1.2,0.0910) (-1.0,0.2062) (-0.8,0.3788) (-0.6,0.6189) (-0.4,0.8588) (-0.2,1.0000) (0.0,0.9958) (0.2,0.8527) (0.4,0.6224) (0.6,0.3854) (0.8,0.2042) (1.0,0.0942) (1.2,0.0356) (1.4,0.0118) (1.6,0.0029) (1.8,0.0011) (2.0,0.0001) (2.2,0.0000) (2.4,0.0000) (2.6,0.0000) (2.8,0.0000) (3.0,0.0000) (3.2,0.0000) (3.4,0.0000) (3.6,0.0000) (3.8,0.0000) (4.0,0.0000) (4.2,0.0000) (4.4,0.0000) (4.6,0.0000) (4.8,0.0000) (5.0,0.0000) (5.2,0.0000) (5.4,0.0000) (5.6,0.0000) (5.8,0.0000) (6.0,0.0000) (6.2,0.0000) (6.4,0.0000) (6.6,0.0000) (6.8,0.0000) (7.0,0.0000) (7.2,0.0000) (7.4,0.0000) (7.6,0.0000) (7.8,0.0000) (8.0,0)} & \histpanel{(-8.0,0.0000) (-7.8,0.0000) (-7.6,0.0000) (-7.4,0.0000) (-7.2,0.0000) (-7.0,0.0000) (-6.8,0.0000) (-6.6,0.0000) (-6.4,0.0000) (-6.2,0.0000) (-6.0,0.0000) (-5.8,0.0000) (-5.6,0.0000) (-5.4,0.0000) (-5.2,0.0000) (-5.0,0.0000) (-4.8,0.0000) (-4.6,0.0000) (-4.4,0.0000) (-4.2,0.0000) (-4.0,0.0000) (-3.8,0.0000) (-3.6,0.0000) (-3.4,0.0000) (-3.2,0.0000) (-3.0,0.0000) (-2.8,0.0000) (-2.6,0.0000) (-2.4,0.0000) (-2.2,0.0000) (-2.0,0.0000) (-1.8,0.0000) (-1.6,0.0000) (-1.4,0.0000) (-1.2,0.0000) (-1.0,0.0000) (-0.8,0.0007) (-0.6,0.0229) (-0.4,0.2741) (-0.2,0.9679) (0.0,1.0000) (0.2,0.3095) (0.4,0.0287) (0.6,0.0010) (0.8,0.0000) (1.0,0.0000) (1.2,0.0000) (1.4,0.0000) (1.6,0.0000) (1.8,0.0000) (2.0,0.0000) (2.2,0.0000) (2.4,0.0000) (2.6,0.0000) (2.8,0.0000) (3.0,0.0000) (3.2,0.0000) (3.4,0.0000) (3.6,0.0000) (3.8,0.0000) (4.0,0.0000) (4.2,0.0000) (4.4,0.0000) (4.6,0.0000) (4.8,0.0000) (5.0,0.0000) (5.2,0.0000) (5.4,0.0000) (5.6,0.0000) (5.8,0.0000) (6.0,0.0000) (6.2,0.0000) (6.4,0.0000) (6.6,0.0000) (6.8,0.0000) (7.0,0.0000) (7.2,0.0000) (7.4,0.0000) (7.6,0.0000) (7.8,0.0000) (8.0,0)} & \histpanel{(-8.0,0.0000) (-7.8,0.0000) (-7.6,0.0000) (-7.4,0.0000) (-7.2,0.0000) (-7.0,0.0000) (-6.8,0.0000) (-6.6,0.0000) (-6.4,0.0000) (-6.2,0.0000) (-6.0,0.0000) (-5.8,0.0000) (-5.6,0.0000) (-5.4,0.0000) (-5.2,0.0000) (-5.0,0.0000) (-4.8,0.0000) (-4.6,0.0000) (-4.4,0.0000) (-4.2,0.0000) (-4.0,0.0000) (-3.8,0.0000) (-3.6,0.0000) (-3.4,0.0000) (-3.2,0.0000) (-3.0,0.0000) (-2.8,0.0000) (-2.6,0.0000) (-2.4,0.0000) (-2.2,0.0000) (-2.0,0.0000) (-1.8,0.0000) (-1.6,0.0000) (-1.4,0.0000) (-1.2,0.0000) (-1.0,0.0000) (-0.8,0.0000) (-0.6,0.0000) (-0.4,0.0016) (-0.2,1.0000) (0.0,0.9333) (0.2,0.0023) (0.4,0.0000) (0.6,0.0000) (0.8,0.0000) (1.0,0.0000) (1.2,0.0000) (1.4,0.0000) (1.6,0.0000) (1.8,0.0000) (2.0,0.0000) (2.2,0.0000) (2.4,0.0000) (2.6,0.0000) (2.8,0.0000) (3.0,0.0000) (3.2,0.0000) (3.4,0.0000) (3.6,0.0000) (3.8,0.0000) (4.0,0.0000) (4.2,0.0000) (4.4,0.0000) (4.6,0.0000) (4.8,0.0000) (5.0,0.0000) (5.2,0.0000) (5.4,0.0000) (5.6,0.0000) (5.8,0.0000) (6.0,0.0000) (6.2,0.0000) (6.4,0.0000) (6.6,0.0000) (6.8,0.0000) (7.0,0.0000) (7.2,0.0000) (7.4,0.0000) (7.6,0.0000) (7.8,0.0000) (8.0,0)} & \histpanel{(-8.0,0.0000) (-7.8,0.0000) (-7.6,0.0000) (-7.4,0.0000) (-7.2,0.0000) (-7.0,0.0000) (-6.8,0.0000) (-6.6,0.0000) (-6.4,0.0000) (-6.2,0.0000) (-6.0,0.0000) (-5.8,0.0000) (-5.6,0.0000) (-5.4,0.0000) (-5.2,0.0000) (-5.0,0.0000) (-4.8,0.0000) (-4.6,0.0000) (-4.4,0.0000) (-4.2,0.0000) (-4.0,0.0000) (-3.8,0.0000) (-3.6,0.0000) (-3.4,0.0000) (-3.2,0.0000) (-3.0,0.0000) (-2.8,0.0000) (-2.6,0.0000) (-2.4,0.0000) (-2.2,0.0000) (-2.0,0.0000) (-1.8,0.0000) (-1.6,0.0000) (-1.4,0.0000) (-1.2,0.0000) (-1.0,0.0000) (-0.8,0.0000) (-0.6,0.0000) (-0.4,0.0000) (-0.2,0.9388) (0.0,1.0000) (0.2,0.0000) (0.4,0.0000) (0.6,0.0000) (0.8,0.0000) (1.0,0.0000) (1.2,0.0000) (1.4,0.0000) (1.6,0.0000) (1.8,0.0000) (2.0,0.0000) (2.2,0.0000) (2.4,0.0000) (2.6,0.0000) (2.8,0.0000) (3.0,0.0000) (3.2,0.0000) (3.4,0.0000) (3.6,0.0000) (3.8,0.0000) (4.0,0.0000) (4.2,0.0000) (4.4,0.0000) (4.6,0.0000) (4.8,0.0000) (5.0,0.0000) (5.2,0.0000) (5.4,0.0000) (5.6,0.0000) (5.8,0.0000) (6.0,0.0000) (6.2,0.0000) (6.4,0.0000) (6.6,0.0000) (6.8,0.0000) (7.0,0.0000) (7.2,0.0000) (7.4,0.0000) (7.6,0.0000) (7.8,0.0000) (8.0,0)} \\
\end{tabular}
\end{center}

Associe cada inicialização à linha de histogramas correspondente:
\textbf{1.} $\sigma_\theta^2 = 1/(4n)$; \quad
\textbf{2.} $\sigma_\theta^2 = 1/n$; \quad
\textbf{3.} $\sigma_\theta^2 = 2/n$; \quad
\textbf{4.} $\sigma_\theta^2 = 4/n$.

\begin{minipage}[t]{0.47\linewidth}
\begin{enumerate}[a)]
    \item 1-S; 2-R; 3-P; 4-Q
    \item 1-R; 2-Q; 3-S; 4-P
    \item 1-Q; 2-R; 3-P; 4-S
    \item 1-P; 2-Q; 3-S; 4-R
\end{enumerate}
\end{minipage}\hfill
\begin{minipage}[t]{0.47\linewidth}
\begin{enumerate}[a)]
\setcounter{enumii}{4}
    \item 1-S; 2-P; 3-R; 4-Q
    \item 1-R; 2-S; 3-Q; 4-P
    \item 1-P; 2-S; 3-Q; 4-R
    \item 1-Q; 2-P; 3-R; 4-S
\end{enumerate}
\end{minipage}

%% ---------- Q2: OBJ-J composição de ReLUs (adaptada de Prince, cap. 3) ----------
\pagebreak
\item (1{,}0) Considere o MLP da figura abaixo, com uma entrada $x$, três unidades ocultas com ativação $\mathrm{a}[z] = \mathrm{ReLU}[z] = \max(0, z)$ e uma saída linear
\[
y = \phi_0 + \phi_1 h_1 + \phi_2 h_2 + \phi_3 h_3, \qquad h_k = \mathrm{a}[\theta_{k0} + \theta_{k1} x],
\]
com $\phi_0 = 0{,}2$. Os gráficos mostram, para $x \in [0, 2]$, as pré-ativações, as ativações, as ativações ponderadas pelos pesos da camada de saída e o viés $\phi_0$ da saída.

\begin{center}
\begin{tikzpicture}[>=stealth, scale=0.85, every node/.style={font=\small}]
    \node[circle, draw, minimum size=8mm] (one) at (0, 1.3) {$1$};
    \node[circle, draw, minimum size=8mm] (x)   at (0, 0)   {$x$};
    \node[circle, draw, minimum size=8mm] (h1)  at (3, 1.3)  {$h_1$};
    \node[circle, draw, minimum size=8mm] (h2)  at (3, 0)    {$h_2$};
    \node[circle, draw, minimum size=8mm] (h3)  at (3, -1.3) {$h_3$};
    \node[circle, draw, minimum size=8mm] (one2) at (3, 2.6) {$1$};
    \node[circle, draw, minimum size=8mm] (y)   at (6, 0.65)   {$y$};
    \draw[->, gray!70] (one) -- (h1); \draw[->, gray!70] (one) -- (h2); \draw[->, gray!70] (one) -- (h3);
    \draw[->] (x) -- (h1); \draw[->] (x) -- (h2); \draw[->] (x) -- (h3);
    \node[font=\scriptsize] at (1.5, 1.55) {$\theta_{k0}$};
    \node[font=\scriptsize] at (1.5, -0.95) {$\theta_{k1}$};
    \draw[->, gray!70] (one2) -- node[above right=-2pt, font=\scriptsize]{$\phi_0$} (y);
    \draw[->] (h1) -- node[above, font=\scriptsize]{$\phi_1$} (y);
    \draw[->] (h2) -- node[above, font=\scriptsize, pos=0.4]{$\phi_2$} (y);
    \draw[->] (h3) -- node[below right=-2pt, font=\scriptsize]{$\phi_3$} (y);
\end{tikzpicture}

\vspace{0.4em}
\setlength{\tabcolsep}{4pt}
\begin{tabular}{cccc}
\fpA{(0,0.6) -- (2,-0.4)}{$\theta_{10} + \theta_{11}x$} &
\fpA{(0,-0.2) -- (2,0.8)}{$\theta_{20} + \theta_{21}x$} &
\fpA{(0,-0.8) -- (2,0.2)}{$\theta_{30} + \theta_{31}x$} & \\[0.5em]
\fpA{(0,0.6) -- (1.2,0) -- (2,0)}{$h_1$} &
\fpA{(0,0) -- (0.4,0) -- (2,0.8)}{$h_2$} &
\fpA{(0,0) -- (1.6,0) -- (2,0.2)}{$h_3$} & \\[0.5em]
\fpA{(0,0.6) -- (1.2,0) -- (2,0)}{$\phi_1 h_1$} &
\fpA{(0,0) -- (0.4,0) -- (2,-0.8)}{$\phi_2 h_2$} &
\fpA{(0,0) -- (1.6,0) -- (2,0.6)}{$\phi_3 h_3$} &
\fpA{(0,0.2) -- (2,0.2)}{$\phi_0$} \\
\end{tabular}
\end{center}

Qual dos gráficos abaixo representa a saída $y$ desse MLP para $x \in [0, 2]$?

\begin{center}
\setlength{\tabcolsep}{1pt}
\begin{tabular}{r@{\,}lr@{\,}lr@{\,}lr@{\,}l}
\textbf{a)} & \fpB{(0,0.6) -- (0.4,0.4) -- (1.2,-0.4) -- (1.6,-0.6) -- (2,-0.2)} &
\textbf{b)} & \fpB{(0,0.6) -- (0.4,0.4) -- (1.2,0.4) -- (1.6,0.6) -- (2,1.0)} &
\textbf{c)} & \fpB{(0,0.8) -- (0.4,0.6) -- (1.2,-0.2) -- (1.6,-0.4) -- (2,0)} &
\textbf{d)} & \fpB{(0,-1.4) -- (2,-0.4)} \\[1.4em]
\textbf{e)} & \fpB{(0,0.6) -- (0.4,0.6) -- (1.2,0.2) -- (1.6,0.2) -- (2,0.8)} &
\textbf{f)} & \fpB{(0,0.8) -- (0.4,0.6) -- (1.2,0.6) -- (1.6,0.8) -- (2,1.2)} &
\textbf{g)} & \fpB{(0,0.8) -- (0.4,0.6) -- (1.2,-0.2) -- (1.6,-0.4) -- (2,-0.6)} &
\textbf{h)} & \fpB{(0,0.2) -- (0.4,0.2) -- (1.2,-0.2) -- (1.6,-0.2) -- (2,0.4)} \\
\end{tabular}
\end{center}

%% ---------- Q3: OBJ-J rede de Perceptrons -> fronteira de decisão em V ----------
\pagebreak
\item (1{,}0) Na rede abaixo, cada círculo é uma unidade que aplica $\mathrm{sinal}$ à soma ponderada das suas entradas somada ao viés indicado \emph{dentro} do círculo, e os números sobre as arestas são os pesos. Considere $\mathrm{sinal}(z) = +1$ para $z \geq 0$ e $\mathrm{sinal}(z) = -1$ para $z < 0$.

\begin{center}
\begin{tikzpicture}[>=stealth]
  \node[mlpin]   (x1) at (0,0.9)   {$x_1$};
  \node[mlpin]   (x2) at (0,-0.9)  {$x_2$};
  \node[mlpunit] (h1) at (2.6,0.9)  {$1$};
  \node[mlpunit] (h2) at (2.6,-0.9) {$-3$};
  \node[mlpunit] (o)  at (5.0,0)    {$-1$};
  \node[font=\scriptsize, above=0pt of h1] {$h_1$};
  \node[font=\scriptsize, below=0pt of h2] {$h_2$};
  \node[mlpin, right=0.6cm of o] (yh) {$\hat{y}$};
  \draw[->] (x1) -- node[mlpw, above]{$-1$} (h1);
  \draw[->] (x2) -- node[mlpw, pos=0.25, above left=-1pt]{$1$} (h1);
  \draw[->] (x1) -- node[mlpw, pos=0.25, below left=-1pt]{$1$} (h2);
  \draw[->] (x2) -- node[mlpw, below]{$1$} (h2);
  \draw[->] (h1) -- node[mlpw, above right=-2pt]{$1$} (o);
  \draw[->] (h2) -- node[mlpw, below right=-2pt]{$1$} (o);
  \draw[->] (o) -- (yh);
\end{tikzpicture}
\end{center}

Em qual dos gráficos abaixo a região sombreada corresponde exatamente ao conjunto de pontos $(x_1, x_2)$ classificados como $\hat{y} = +1$?

\begin{center}
\setlength{\tabcolsep}{1pt}
\begin{tabular}{r@{\,}lr@{\,}lr@{\,}lr@{\,}l}
\textbf{a)} & \regpanel{(0,0) -- (1,0) -- (2,1) -- (3,0) -- (4,0) -- (4,4) -- (0,4) -- cycle} &
\textbf{b)} & \regpanel{(0,0) -- (4,0) -- (4,3) -- (2,1) -- (0,3) -- cycle} &
\textbf{c)} & \regpanel{(0,3) -- (1,2) -- (3,4) -- (0,4) -- cycle} &
\textbf{d)} & \regpanel{(2,1) -- (4,3) -- (4,0) -- (3,0) -- cycle} \\[0.5em]
\textbf{e)} & \regpanel{(2,1) -- (0,3) -- (0,0) -- (1,0) -- cycle} &
\textbf{f)} & \regpanel{(0,0) -- (1,0) -- (4,3) -- (4,4) -- (0,4) -- cycle} &
\textbf{g)} & \regpanel{(0,3) -- (2,1) -- (4,3) -- (4,4) -- (0,4) -- cycle} &
\textbf{h)} & \regpanel{(1,0) -- (2,1) -- (3,0) -- cycle} \\
\end{tabular}
\end{center}

%% ---------- Q4: softmax + entropia cruzada, deslocamento e escala dos logits ----------
\item (1{,}0) Um classificador com três classes produz os \textit{logits} $\mathbf{z} = \begin{bmatrix} 1 & 3 & 1 \end{bmatrix}$ para uma instância cuja classe correta é a primeira. A saída passa por uma SoftMax e o custo é a entropia cruzada. Considere também duas modificações dos \textit{logits}: \textbf{(A)} somar $5$ a todos os componentes, obtendo $\begin{bmatrix} 6 & 8 & 6 \end{bmatrix}$, e \textbf{(B)} multiplicar todos os componentes por $2$, obtendo $\begin{bmatrix} 2 & 6 & 2 \end{bmatrix}$. Use $e^{2} \approx 7{,}39$, $e^{4} \approx 54{,}60$, $\ln 9{,}39 \approx 2{,}24$ e $\ln 56{,}60 \approx 4{,}04$. Quais são os valores da entropia cruzada para os \textit{logits} originais, para (A) e para (B), nessa ordem?

\begin{minipage}[t]{0.47\linewidth}
\begin{enumerate}[a)]
    \item $2{,}24$; $2{,}24$; $4{,}04$
    \item $2{,}24$; $7{,}24$; $4{,}04$
    \item $2{,}24$; $2{,}24$; $2{,}24$
    \item $2{,}24$; $2{,}24$; $4{,}48$
\end{enumerate}
\end{minipage}\hfill
\begin{minipage}[t]{0.47\linewidth}
\begin{enumerate}[a)]
\setcounter{enumii}{4}
    \item $2{,}24$; $7{,}24$; $4{,}48$
    \item $2{,}24$; $7{,}24$; $2{,}24$
    \item $2{,}24$; $2{,}24$; $1{,}12$
    \item $2{,}24$; $7{,}24$; $1{,}12$
\end{enumerate}
\end{minipage}

%% ---------- Q5: forward pass vetorizado com dropout ----------
\pagebreak
\item (1{,}0) Durante o treinamento, uma camada de \textit{dropout} com taxa $p$ sorteia, a cada iteração, uma máscara binária com o mesmo formato da sua entrada: cada posição vale $0$ com probabilidade $p$ e $1$ caso contrário. A camada multiplica a entrada, elemento a elemento, pela máscara e multiplica o resultado por $1/(1-p)$, para que o valor esperado de cada ativação não mude (\textit{inverted dropout}).

Considere o MLP $1 \times 2 \times 2 \times 1$ abaixo, com $\mathbf{z}^{(l)} = \mathbf{a}^{(l-1)}{\mathbf{\Theta}^{(l)}}^{\top} + \mathbf{b}^{(l)}$, ReLU nas duas camadas ocultas, saída linear e uma camada de \textit{dropout} com $p = 0{,}5$ aplicada sobre $\mathbf{a}^{(1)}$:

\begin{center}
\begin{tikzpicture}[>=stealth, yscale=0.8, every node/.style={font=\small}]
    % entrada
    \node[circle, draw] (x) at (0, 0) {$x$};
    % primeira camada oculta (ReLU)
    \node[circle, draw] (a1) at (2.4, 1.1)  {$a^{(1)}_1$};
    \node[circle, draw] (a2) at (2.4, -1.1) {$a^{(1)}_2$};
    % dropout: um quadrado por unidade
    \node[rectangle, draw, fill=black!8, minimum size=7mm] (d1) at (4.2, 1.1)  {};
    \node[rectangle, draw, fill=black!8, minimum size=7mm] (d2) at (4.2, -1.1) {};
    % segunda camada oculta (ReLU)
    \node[circle, draw] (c1) at (6.6, 1.1)  {$a^{(2)}_1$};
    \node[circle, draw] (c2) at (6.6, -1.1) {$a^{(2)}_2$};
    % saída linear
    \node[circle, draw] (y) at (9.0, 0) {$\hat{y}$};
    \draw[->] (x) -- (a1); \draw[->] (x) -- (a2);
    \draw[->] (a1) -- (d1); \draw[->] (a2) -- (d2);
    \draw[->] (d1) -- (c1); \draw[->] (d1) -- (c2);
    \draw[->] (d2) -- (c1); \draw[->] (d2) -- (c2);
    \draw[->] (c1) -- (y); \draw[->] (c2) -- (y);
    % rótulos das camadas
    \node[font=\footnotesize] at (1.2, 2.0)  {$\mathbf{\Theta}^{(1)}, \mathbf{b}^{(1)}$};
    \node[font=\footnotesize] at (5.4, 2.0)  {$\mathbf{\Theta}^{(2)}, \mathbf{b}^{(2)}$};
    \node[font=\footnotesize] at (8.6, 1.1) {$\mathbf{\Theta}^{(3)}, b^{(3)}$};
    \node[font=\footnotesize, below=0.12cm of a2] {ReLU};
    \node[font=\footnotesize, below=0.12cm of d2] {\textit{dropout}};
    \node[font=\footnotesize, below=0.12cm of c2] {ReLU};
    \node[font=\footnotesize, below=0.12cm of y]  {linear};
\end{tikzpicture}
\end{center}

Os parâmetros, a instância de treino e a máscara sorteada nesta iteração são
\[
x = 2, \quad
\mathbf{\Theta}^{(1)} = \begin{bmatrix} -1 \\ 2 \end{bmatrix}, \quad
\mathbf{b}^{(1)} = \begin{bmatrix} 3 & 1 \end{bmatrix}, \quad
\mathbf{\Theta}^{(2)} = \begin{bmatrix} 3 & 0{,}5 \\ -2 & -2 \end{bmatrix}, \quad
\mathbf{b}^{(2)} = \begin{bmatrix} 1 & 0 \end{bmatrix},
\]
\[
\mathbf{\Theta}^{(3)} = \begin{bmatrix} -1 & -3 \end{bmatrix}, \quad
b^{(3)} = 1, \quad
p = 0{,}5, \quad
\text{máscara} = \begin{bmatrix} 1 & 0 \end{bmatrix}.
\]
Qual é a saída $\hat{y}$ da rede nesta iteração de treinamento?

\begin{center}
\begin{tabular}{@{}l@{\hspace{3em}}l@{\hspace{3em}}l@{\hspace{3em}}l@{}}
a) $\hat{y} = -5{,}5$ & b) $\hat{y} = -6$ & c) $\hat{y} = -3$ & d) $\hat{y} = -1{,}5$ \\[0.3em]
e) $\hat{y} = -5$ & f) $\hat{y} = 6$ & g) $\hat{y} = 0$ & h) $\hat{y} = -12$ \\
\end{tabular}
\end{center}

%% ---------- Q6: OBJ-H otimizadores -> comportamento observado ----------
\item (1{,}0) Associe cada otimizador (coluna da esquerda) ao comportamento observado durante o treinamento (coluna da direita).

\begin{center}
\begin{tabular}{p{0.22\linewidth} p{0.70\linewidth}}
\textbf{1.} \textit{Momentum} & \textbf{(i)} a taxa efetiva de cada parâmetro volta a crescer quando os gradientes recentes desse parâmetro ficam pequenos. \\
\textbf{2.} Nesterov & \textbf{(ii)} em um vale estreito, as oscilações transversais se cancelam e o avanço ao longo do vale acelera. \\
\textbf{3.} AdaGrad & \textbf{(iii)} a taxa efetiva de cada parâmetro só diminui, e após muitas iterações o treinamento praticamente estaciona. \\
\textbf{4.} RMSProp & \textbf{(iv)} o gradiente é avaliado no ponto para onde a memória levaria, e o passo é corrigido antes de uma curva do vale. \\
\end{tabular}
\end{center}

\begin{minipage}[t]{0.47\linewidth}
\begin{enumerate}[a)]
    \item 1-(iv); 2-(ii); 3-(iii); 4-(i)
    \item 1-(iii); 2-(i); 3-(ii); 4-(iv)
    \item 1-(ii); 2-(iv); 3-(i); 4-(iii)
    \item 1-(i); 2-(iii); 3-(ii); 4-(iv)
\end{enumerate}
\end{minipage}\hfill
\begin{minipage}[t]{0.47\linewidth}
\begin{enumerate}[a)]
\setcounter{enumii}{4}
    \item 1-(iv); 2-(ii); 3-(i); 4-(iii)
    \item 1-(iii); 2-(i); 3-(iv); 4-(ii)
    \item 1-(i); 2-(iii); 3-(iv); 4-(ii)
    \item 1-(ii); 2-(iv); 3-(iii); 4-(i)
\end{enumerate}
\end{minipage}

%% ---------- Q7: dois objetivos, direção que melhora L1 sem alterar L2 ----------
\pagebreak
\item (1{,}0) Uma rede é treinada com dois objetivos, representados pelas funções de custo $\mathcal{L}_1$ e $\mathcal{L}_2$. No ponto atual dos parâmetros, os gradientes são os mostrados abaixo.

\begin{center}
\begin{minipage}{0.3\linewidth}
\[
\nabla_{\boldsymbol{\theta}}\mathcal{L}_1 = \begin{bmatrix} 4 \\ 1 \end{bmatrix}
\]
\[
\nabla_{\boldsymbol{\theta}}\mathcal{L}_2 = \begin{bmatrix} -2 \\ 1 \end{bmatrix}
\]
\end{minipage}%
\begin{minipage}{0.68\linewidth}\centering
\begin{tikzpicture}[scale=0.98, >=stealth]
  \foreach \g in {-3,-2,-1,1,2,3,4} { \draw[gray!20] (\g,-2.5) -- (\g,2.5); }
  \foreach \g in {-2,-1,1,2} { \draw[gray!20] (-3.5,\g) -- (4.5,\g); }
  \draw[gray!60, ->] (-3.5,0) -- (4.7,0) node[right, font=\small, black] {$\theta_1$};
  \draw[gray!60, ->] (0,-2.5) -- (0,2.7) node[above, font=\small, black] {$\theta_2$};
  \foreach \t in {-3,-2,-1,1,2,3,4} \node[font=\scriptsize, below, gray!60!black, inner sep=1.5pt] at (\t,0) {$\t$};
  \foreach \t in {-2,-1,1,2} \node[font=\scriptsize, left, gray!60!black, inner sep=1.5pt] at (0,\t) {$\t$};
  \draw[->, very thick] (0,0) -- (4,1) node[above right, font=\small] {$\nabla_{\boldsymbol{\theta}}\mathcal{L}_1$};
  \draw[->, very thick, black!55] (0,0) -- (-2,1) node[above left, font=\small, black] {$\nabla_{\boldsymbol{\theta}}\mathcal{L}_2$};
\end{tikzpicture}
\end{minipage}
\end{center}

Você percebe que o passo usual do SGD sobre $\mathcal{L}_1$ melhora o objetivo 1, mas degrada o objetivo 2.

Considere atualizações da forma $\boldsymbol{\theta}_{t+1} = \boldsymbol{\theta}_t + \eta\,\Delta\boldsymbol{\theta}$, com $\eta > 0$ pequeno. Pela aproximação de primeira ordem, a variação de cada função de custo é
\[
\mathcal{L}_k(\boldsymbol{\theta}_t + \eta\,\Delta\boldsymbol{\theta}) - \mathcal{L}_k(\boldsymbol{\theta}_t)
\approx \eta\,\Delta\boldsymbol{\theta}^{\top}\nabla_{\boldsymbol{\theta}}\mathcal{L}_k
= \eta\,\|\Delta\boldsymbol{\theta}\|\,\|\nabla_{\boldsymbol{\theta}}\mathcal{L}_k\|\cos(\alpha_k),
\]
em que $\alpha_k$ é o ângulo entre $\Delta\boldsymbol{\theta}$ e $\nabla_{\boldsymbol{\theta}}\mathcal{L}_k$. Qual direção $\Delta\boldsymbol{\theta}$ reduz $\mathcal{L}_1$ e mantém $\mathcal{L}_2$ inalterada segundo essa aproximação?

\begin{center}
\begin{tabular}{@{}llll@{}}
a) $\Delta\boldsymbol{\theta} = \bigl[\,-4 \quad -1\,\bigr]^{\top}$ & b) $\Delta\boldsymbol{\theta} = \bigl[\,-1 \quad -1\,\bigr]^{\top}$ & c) $\Delta\boldsymbol{\theta} = \bigl[\,-2 \quad -2\,\bigr]^{\top}$ & d) $\Delta\boldsymbol{\theta} = \bigl[\,-1 \quad -2\,\bigr]^{\top}$ \\[0.3em]
e) $\Delta\boldsymbol{\theta} = \bigl[\,-2 \quad 1\,\bigr]^{\top}$ & f) $\Delta\boldsymbol{\theta} = \bigl[\,1 \quad 2\,\bigr]^{\top}$ & g) $\Delta\boldsymbol{\theta} = \bigl[\,1 \quad -4\,\bigr]^{\top}$ & h) $\Delta\boldsymbol{\theta} = \bigl[\,-6 \quad 0\,\bigr]^{\top}$ \\
\end{tabular}
\end{center}

%% ---------- Q8: Adam e Muon a partir da SVD fornecida do momentum ----------
\item (1{,}0) Em uma camada com matriz de pesos $\mathbf{W}$ de formato $2 \times 2$, o \textit{momentum} acumulado nesta iteração e a sua SVD, $\mathbf{M}_{t+1} = \mathbf{U}\boldsymbol{\Sigma}\mathbf{V}^{\top}$, são
\[
\mathbf{M}_{t+1} = \begin{bmatrix} 12 & 4 \\ 11 & 12 \end{bmatrix}, \qquad
\mathbf{U} = \frac{1}{5}\begin{bmatrix} 3 & -4 \\ 4 & 3 \end{bmatrix}, \qquad
\boldsymbol{\Sigma} = \begin{bmatrix} 20 & 0 \\ 0 & 5 \end{bmatrix}, \qquad
\mathbf{V} = \frac{1}{5}\begin{bmatrix} 4 & -3 \\ 3 & 4 \end{bmatrix}.
\]
Os dois otimizadores atualizam a camada como $\mathbf{W}_{t+1} = \mathbf{W}_t - \eta\,\mathbf{D}$, cada um com a sua matriz de direção $\mathbf{D}$. Suponha que o Adam tenha histórico de gradientes estável, com $\widetilde{\mathbf{m}}_{t+1} = \mathbf{M}_{t+1}$ e $\widetilde{\mathbf{v}}_{t+1} = \widetilde{\mathbf{m}}_{t+1}^{2}$ (entrada a entrada, $\epsilon \approx 0$), e que o Muon receba o mesmo $\mathbf{M}_{t+1}$. Quais são $\mathbf{D}_{\text{Adam}}$ e $\mathbf{D}_{\text{Muon}}$?

\begin{center}
{\renewcommand{\arraystretch}{1.0}\small
\begin{tabular}{@{}l@{\hspace{1em}}l@{}}
a) $\mathbf{D}_{\text{Adam}} = \begin{bmatrix} 1 & 1 \\ 1 & 1 \end{bmatrix}$, \ $\mathbf{D}_{\text{Muon}} = \begin{bmatrix} 0{,}96 & 0{,}28 \\ -0{,}28 & 0{,}96 \end{bmatrix}$ & e) $\mathbf{D}_{\text{Adam}} = \begin{bmatrix} 0{,}96 & -0{,}28 \\ 0{,}28 & 0{,}96 \end{bmatrix}$, \ $\mathbf{D}_{\text{Muon}} = \begin{bmatrix} 0{,}96 & 0{,}28 \\ -0{,}28 & 0{,}96 \end{bmatrix}$ \\[0.9em]
b) $\mathbf{D}_{\text{Adam}} = \begin{bmatrix} 0{,}96 & -0{,}28 \\ 0{,}28 & 0{,}96 \end{bmatrix}$, \ $\mathbf{D}_{\text{Muon}} = \begin{bmatrix} 0{,}96 & -0{,}28 \\ 0{,}28 & 0{,}96 \end{bmatrix}$ & f) $\mathbf{D}_{\text{Adam}} = \begin{bmatrix} 1 & 1 \\ 1 & 1 \end{bmatrix}$, \ $\mathbf{D}_{\text{Muon}} = \begin{bmatrix} 0{,}96 & -0{,}28 \\ 0{,}28 & 0{,}96 \end{bmatrix}$ \\[0.9em]
c) $\mathbf{D}_{\text{Adam}} = \begin{bmatrix} 1 & 1 \\ 1 & 1 \end{bmatrix}$, \ $\mathbf{D}_{\text{Muon}} = \begin{bmatrix} 1 & 1 \\ 1 & 1 \end{bmatrix}$ & g) $\mathbf{D}_{\text{Adam}} = \begin{bmatrix} 1 & 1 \\ 1 & 1 \end{bmatrix}$, \ $\mathbf{D}_{\text{Muon}} = \begin{bmatrix} 0{,}6 & 0{,}2 \\ 0{,}55 & 0{,}6 \end{bmatrix}$ \\[0.9em]
d) $\mathbf{D}_{\text{Adam}} = \begin{bmatrix} 0{,}96 & -0{,}28 \\ 0{,}28 & 0{,}96 \end{bmatrix}$, \ $\mathbf{D}_{\text{Muon}} = \begin{bmatrix} 0{,}6 & 0{,}2 \\ 0{,}55 & 0{,}6 \end{bmatrix}$ & h) $\mathbf{D}_{\text{Adam}} = \begin{bmatrix} 0{,}96 & -0{,}28 \\ 0{,}28 & 0{,}96 \end{bmatrix}$, \ $\mathbf{D}_{\text{Muon}} = \begin{bmatrix} 1 & 1 \\ 1 & 1 \end{bmatrix}$ \\
\end{tabular}}
\end{center}

%% ---------- Q9: backward pass com sigmoide + entropia cruzada binária ----------
\pagebreak
\item (1{,}0) Considere o MLP abaixo, com duas unidades ReLU na camada oculta e uma unidade sigmoide na saída, treinado com a entropia cruzada binária $J = -\bigl(y\ln\hat{y} + (1-y)\ln(1-\hat{y})\bigr)$. A entrada é $x = 2$, o rótulo é $y = 1$ e os parâmetros são
\[
\mathbf{\Theta}^{(1)} = \begin{bmatrix} 1 \\ -0{,}5 \end{bmatrix}, \quad
\mathbf{b}^{(1)} = \begin{bmatrix} -1 & 2 \end{bmatrix}, \quad
\mathbf{\Theta}^{(2)} = \begin{bmatrix} 2 & -1{,}5 \end{bmatrix}, \quad
b^{(2)} = -0{,}5 .
\]

\begin{center}
\begin{tikzpicture}[>=stealth, every node/.style={font=\small}]
    \node[circle, draw] (x)  at (0, 0)      {$x$};
    \node[circle, draw] (h1) at (2.8, 1.15)  {$a^{(1)}_1$};
    \node[circle, draw] (h2) at (2.8, -1.15) {$a^{(1)}_2$};
    \node[circle, draw] (y)  at (5.6, 0)    {$\hat{y}$};
    \draw[->] (x)  -- node[above left=-0.08cm, font=\footnotesize]{$\theta^{(1)}_{11}$} (h1);
    \draw[->] (x)  -- node[below left=-0.08cm, font=\footnotesize]{$\theta^{(1)}_{21}$} (h2);
    \draw[->] (h1) -- node[above right=-0.08cm, font=\footnotesize]{$\theta^{(2)}_{11}$} (y);
    \draw[->] (h2) -- node[below right=-0.08cm, font=\footnotesize]{$\theta^{(2)}_{12}$} (y);
    \draw[->] (2.8, 2.35)  -- node[right, font=\footnotesize]{$b^{(1)}_1$} (h1);
    \draw[->] (2.8, -2.35) -- node[right, font=\footnotesize]{$b^{(1)}_2$} (h2);
    \draw[->] (5.6, 1.15)  -- node[right, font=\footnotesize]{$b^{(2)}$}   (y);
    \node[left=0.10cm of h1, font=\footnotesize] {ReLU};
    \node[left=0.10cm of h2, font=\footnotesize] {ReLU};
    \node[below=0.30cm of y, font=\footnotesize] {sigmoide};
\end{tikzpicture}
\end{center}

Sabendo que $\sigma(0) = 0{,}5$ e que a derivada da entropia cruzada binária em relação à saída é
\[
\frac{\partial J}{\partial \hat{y}} = -\frac{y}{\hat{y}} + \frac{1 - y}{1 - \hat{y}},
\]
qual o valor de $\dfrac{\partial J}{\partial \theta^{(1)}_{21}}$?

\begin{minipage}[t]{0.47\linewidth}
\begin{enumerate}[a)]
    \item $+0{,}75$
    \item $-1{,}0$
    \item $-1{,}5$
    \item $+0{,}375$
\end{enumerate}
\end{minipage}\hfill
\begin{minipage}[t]{0.47\linewidth}
\begin{enumerate}[a)]
\setcounter{enumii}{4}
    \item $+1{,}5$
    \item $+6{,}0$
    \item $-2{,}0$
    \item $-0{,}375$
\end{enumerate}
\end{minipage}

%% ---------- Q10: receita de máxima verossimilhança com Poisson ----------
\item (1{,}0) Você quer prever o número de chamadas que uma central telefônica recebe por hora, $y \in \{0, 1, 2, \ldots\}$, modelando a saída com a distribuição de Poisson
\[
\Pr(y \mid \lambda) = \frac{\lambda^{y}\,e^{-\lambda}}{y!}, \qquad \lambda > 0 ,
\]
em que a rede prediz $\lambda^{(i)} = f(\mathbf{x}^{(i)}; \boldsymbol{\theta})$ para cada uma das $N$ observações independentes. Seguindo a receita de máxima verossimilhança vista em aula e descartando tudo que não altera o ponto de mínimo em $\boldsymbol{\theta}$, qual combinação de função de custo e função de ativação na saída é adequada? As ativações consideradas nas alternativas são
\[
\sigma(z) = \frac{1}{1 + e^{-z}}, \qquad
\mathrm{softplus}(z) = \ln\bigl(1 + e^{z}\bigr), \qquad
\text{identidade: } f(z) = z, \qquad
\mathrm{ReLU}(z) = \max(0, z).
\]

\begin{enumerate}[a)]
    \item $\mathcal{L} = \sum_{i}\bigl[\lambda^{(i)} - y^{(i)}\ln\lambda^{(i)}\bigr]$, com sigmoide na camada de saída.
    \item $\mathcal{L} = \sum_{i}\bigl[\lambda^{(i)} - y^{(i)}\ln\lambda^{(i)}\bigr]$, com \textit{softplus} na camada de saída.
    \item $\mathcal{L} = \sum_{i}\bigl[\lambda^{(i)} - y^{(i)}\ln\lambda^{(i)}\bigr]$, com identidade (linear) na camada de saída.
    \item $\mathcal{L} = \sum_{i}\bigl[\lambda^{(i)} - y^{(i)}\ln\lambda^{(i)}\bigr]$, com ReLU na camada de saída.
    \item $\mathcal{L} = \sum_{i}\bigl[y^{(i)}\ln\lambda^{(i)} - \lambda^{(i)}\bigr]$, com \textit{softplus} na camada de saída.
    \item $\mathcal{L} = \sum_{i}\bigl[y^{(i)}\ln\lambda^{(i)} - \lambda^{(i)}\bigr]$, com sigmoide na camada de saída.
    \item $\mathcal{L} = \sum_{i}\bigl[y^{(i)}\ln\lambda^{(i)} - \lambda^{(i)}\bigr]$, com ReLU na camada de saída.
    \item $\mathcal{L} = \sum_{i}\bigl[y^{(i)}\ln\lambda^{(i)} - \lambda^{(i)}\bigr]$, com identidade (linear) na camada de saída.
\end{enumerate}

\end{enumerate}

\end{document}
